% =====================================================================
%  INDRA 2.0 — Technical Report
%  Versions 3.8 and 3.9
%  Compile:  pdflatex INDRA_Technical_Report.tex   (run twice for TOC)
% =====================================================================
\documentclass[11pt,a4paper]{article}

\usepackage[T1]{fontenc}
\usepackage[utf8]{inputenc}
\usepackage{lmodern}
\usepackage[margin=2.4cm,top=2.6cm,bottom=2.6cm]{geometry}
\usepackage{amsmath,amssymb}
\usepackage{booktabs}
\usepackage{array}
\usepackage{longtable}
\usepackage{xcolor}
\usepackage{graphicx}
\usepackage{float}
\usepackage{enumitem}
\usepackage{titlesec}
\usepackage{fancyhdr}
\usepackage[most]{tcolorbox}
\usepackage{microtype}
\usepackage[hidelinks]{hyperref}

% ---------------------------------------------------------------- colours
\definecolor{navy}{HTML}{1E2761}
\definecolor{navylight}{HTML}{3A4A8F}
\definecolor{ice}{HTML}{EEF3FC}
\definecolor{warnred}{HTML}{A8322A}
\definecolor{warnbg}{HTML}{FBEDEB}
\definecolor{goodgreen}{HTML}{1F6A50}
\definecolor{goodbg}{HTML}{E8F3EE}
\definecolor{goldc}{HTML}{9A7524}
\definecolor{goldbg}{HTML}{FBF4E4}
\definecolor{greytext}{HTML}{555D68}
\definecolor{rulegrey}{HTML}{C9D2E4}

% ---------------------------------------------------------------- headings
\titleformat{\section}
  {\color{navy}\normalfont\Large\bfseries}{\thesection}{0.7em}{}
\titleformat{\subsection}
  {\color{navylight}\normalfont\large\bfseries}{\thesubsection}{0.6em}{}
\titleformat{\subsubsection}
  {\color{navylight}\normalfont\normalsize\bfseries}{\thesubsubsection}{0.6em}{}
\titlespacing*{\section}{0pt}{1.6em}{0.7em}
\titlespacing*{\subsection}{0pt}{1.1em}{0.5em}

% ---------------------------------------------------------------- headers
\setlength{\headheight}{22pt}
\pagestyle{fancy}
\fancyhf{}
\renewcommand{\headrulewidth}{0.4pt}
\renewcommand{\headrule}{\hbox to\headwidth{\color{rulegrey}\leaders\hrule height \headrulewidth\hfill}}
\fancyhead[L]{\color{greytext}\small INDRA 2.0 --- Technical Report}
\fancyhead[R]{\color{greytext}\small v3.8 \& v3.9}
\fancyfoot[C]{\color{greytext}\small\thepage}

% ---------------------------------------------------------------- boxes
\newtcolorbox{keybox}[1][]{
  breakable, colback=ice, colframe=navy, boxrule=0pt, leftrule=3pt,
  arc=1pt, left=10pt, right=10pt, top=8pt, bottom=8pt, #1}
\newtcolorbox{warnbox}[1][]{
  breakable, colback=warnbg, colframe=warnred, boxrule=0pt, leftrule=3pt,
  arc=1pt, left=10pt, right=10pt, top=8pt, bottom=8pt, #1}
\newtcolorbox{goodbox}[1][]{
  breakable, colback=goodbg, colframe=goodgreen, boxrule=0pt, leftrule=3pt,
  arc=1pt, left=10pt, right=10pt, top=8pt, bottom=8pt, #1}
\newtcolorbox{goldbox}[1][]{
  breakable, colback=goldbg, colframe=goldc, boxrule=0pt, leftrule=3pt,
  arc=1pt, left=10pt, right=10pt, top=8pt, bottom=8pt, #1}

% ---------------------------------------------------------------- macros
\newcommand{\um}{\,$\mu$m}
\newcommand{\code}[1]{\texttt{\small #1}}
\newcommand{\htwoo}{H$_2$O}
\newcommand{\cotwo}{CO$_2$}
\newcommand{\chthreeoh}{CH$_3$OH}
\newcommand{\nhfour}{NH$_4^+$}
\newcommand{\ocn}{OCN$^-$}
\newcommand{\tauv}{$\tau$}

\setlist[itemize]{itemsep=2pt, topsep=4pt, parsep=0pt}
\setlist[enumerate]{itemsep=2pt, topsep=4pt, parsep=0pt}
\newcolumntype{L}[1]{>{\raggedright\arraybackslash}p{#1}}
\renewcommand{\arraystretch}{1.25}

% =====================================================================
\begin{document}

% ---------------------------------------------------------------- title
\begin{titlepage}
\centering
\vspace*{3.2cm}

{\color{navy}\Huge\bfseries INDRA 2.0\par}
\vspace{0.7cm}
{\color{navylight}\Large\itshape Ice Component Identification using\\[2pt]
Deep Learning on Astrophysics Spectra\par}

\vspace{1.4cm}
{\color{rulegrey}\rule{0.45\textwidth}{1pt}}
\vspace{1.2cm}

{\Large\bfseries Complete Technical Report\par}
\vspace{0.35cm}
{\large Pipeline Versions 3.8 and 3.9\par}

\vspace{2.2cm}

{\large School of Physical Sciences\\[2pt]
National Institute of Science Education and Research, Bhubaneswar\par}

\vspace{0.9cm}
{\color{greytext}\normalsize In collaboration with Prathap Rayalacheruvu\\
and Prof.\ Liton Majumdar (SEPS, NISER)\par}

\vfill
{\color{greytext}November 2026}
\end{titlepage}

% ---------------------------------------------------------------- TOC
{\small\setlength{\parskip}{0pt}\tableofcontents}
\thispagestyle{fancy}
\newpage

% =====================================================================
\section{The Problem, Stated Plainly}
% =====================================================================

This report describes a machine learning system built to solve a specific
measurement problem in astronomy. The problem can be understood completely
without any background in chemistry or astronomy, and this opening section
presents it in two ways: first as a physical story, then stripped down to
the mathematical problem that the software actually solves.

\subsection{The physical story}

When a new star forms, it does so inside a cold, dense cloud of gas and
dust. Long before the star ignites, molecules in that cloud freeze onto
microscopic dust grains and build up layers of ice --- not only water ice,
but mixed ices containing carbon dioxide, methanol, formic acid, ammonia
and dozens of other compounds. These ice layers are the chemical raw
material from which planets and comets are later assembled, so knowing what
they are made of is one of the central questions in the study of star
formation.

The James Webb Space Telescope (JWST) can measure these ices indirectly.
Its Mid-Infrared Instrument (MIRI) records how much light the young star
emits at each wavelength between 5 and 28 micrometres. When light passes
through the surrounding ice, each type of molecule absorbs light at its own
characteristic wavelengths. The recorded spectrum therefore looks like a
smooth hill with a series of dips carved into it, and each dip is the
fingerprint of a particular molecule.

If each molecule produced a single clean, isolated dip, reading the
composition off the spectrum would be straightforward. It does not. The
dips overlap heavily. Water absorbs broadly around 3 and 6\um. Methanol and
formic acid both absorb near 9--10\um. Ammonium and several other
nitrogen-bearing species all crowd around 6.8\um. On top of this, silicate
dust --- ordinary rocky mineral grains, not ice at all --- produces an
enormous broad absorption centred at 9.7\um{} that sits underneath
everything else and must be removed before the ice features can even be
seen.

\subsection{How the problem was solved before}

The established approach is manual. An expert astronomer inspects the
spectrum, decides from experience which ices are plausibly present, selects
a corresponding set of laboratory reference spectra --- real ice mixtures
measured in a laboratory under controlled conditions --- and fits a weighted
sum of those references to the observation. On our benchmark source this
produced an excellent result, but it carries three costs:

\begin{itemize}
  \item it takes an expert weeks of work per source;
  \item it does not transfer to a new source without repeating the whole
        manual process;
  \item it is not reproducible --- two astronomers may select different
        reference files and reach different conclusions.
\end{itemize}

\subsection{The computational problem, stripped of astronomy}

Underneath the astronomy, the task is a standard problem in applied
mathematics. The essential statement is this:

\begin{keybox}
\textbf{Problem statement.} We are given an observed signal
$\mathbf{b}\in\mathbb{R}^{3000}$ (the measured absorption spectrum, sampled
on a fixed wavelength grid) and a dictionary
$\mathbf{A}\in\mathbb{R}^{3000\times 1209}$ whose columns are laboratory
reference spectra. We must find a coefficient vector
$\mathbf{x}\in\mathbb{R}^{1209}$ such that
\[
   \mathbf{A}\mathbf{x} \;\approx\; \mathbf{b},
   \qquad \mathbf{x} \ge 0,
   \qquad \|\mathbf{x}\|_0 \ \text{small.}
\]
The non-negativity constraint is physical: a mixture cannot contain a
negative amount of a substance. The sparsity requirement is also physical:
a real ice mantle contains on the order of 10--25 distinguishable
components, not 1209.
\end{keybox}

Three properties make this harder than it first appears.

\begin{description}[style=nextline, leftmargin=0pt, labelindent=0pt, labelwidth=0pt, labelsep=0pt, itemsep=4pt]
  \item[The dictionary is highly correlated.]
  Many columns of $\mathbf{A}$ are nearly identical to one another. The
  same ice mixture measured at 20\,K and at 30\,K produces two spectra that
  differ by less than the measurement noise. Classical least squares
  becomes unstable when the design matrix has near-duplicate columns ---
  the solver can distribute weight between them almost arbitrarily.

  \item[The dictionary is much wider than it is informative.]
  With 1209 candidate columns and an effective resolution that smooths
  away fine detail, many different coefficient vectors reproduce
  $\mathbf{b}$ about equally well. The problem is underdetermined in
  practice even though $3000 > 1209$ formally.

  \item[Ground truth is almost nonexistent.]
  There is exactly one source for which a trusted, published answer
  exists. Everything else must be validated indirectly.
\end{description}

\subsection{What INDRA 2.0 adds}

INDRA 2.0 replaces the manual file-selection step with a trained
classifier. Rather than an astronomer saying \emph{``I believe formic acid
is present, so I will include the formic acid references''}, a machine
learning model examines the spectrum and reports \emph{``given this
absorption pattern, the probability that formic acid is present is
44\%''}. Those probabilities then restrict which columns of $\mathbf{A}$
are offered to the fitting stage.

The system therefore has two distinct stages, and keeping them separate is
the central architectural decision of the project:

\begin{keybox}
\textbf{Stage 1 --- Detection.} Machine learning classifiers read the
spectrum and output a presence probability for each candidate species.
This answers \emph{which} substances are present.

\vspace{4pt}
\textbf{Stage 2 --- Quantification.} A constrained least-squares solver
fits the selected laboratory spectra to the observation. This answers
\emph{how much} of each is present.
\end{keybox}

\subsection{Sources used in this work}

\begin{table}[H]
\centering
\small
\begin{tabular}{@{}L{3.4cm}L{4.6cm}L{6.0cm}@{}}
\toprule
\textbf{Source} & \textbf{Role} & \textbf{Notes} \\
\midrule
NGC 1333 IRAS 2A & Validation benchmark &
Published decomposition available (Rayalacheruvu et al.\ 2025): 76
laboratory spectra, 15 mixture components, 17 chemical species. \\
L1448 IRS2 & Blind generalisation test &
Class 0 protostar, $L\approx5.2\,L_\odot$. No published ground truth. \\
L1448 IRS3B & Blind generalisation test &
Class 0 binary system. No published ground truth. \\
\bottomrule
\end{tabular}
\caption{The three target sources. All three lie in the Perseus molecular
cloud at roughly 300\,pc.}
\end{table}

The distinction between the benchmark and the two blind sources matters
throughout this report. IRAS 2A is the only source against which accuracy
can be measured. L1448 IRS2 and IRS3B are the only sources that test
whether the pipeline generalises --- and, as Section~\ref{sec:leakage}
describes, the difference between these two situations is precisely where
the project's most serious methodological error occurred.

% =====================================================================
\section{Input Data and Preprocessing}
% =====================================================================

\subsection{The laboratory reference library}

The Leiden Ice Database (LIDA) provides 1209 laboratory ice spectra. Each
file is a two-column ASCII table: wavelength (or wavenumber) in the first
column, absorption in the second. The files cover many different mixtures
--- for example ``water:methanol:CO$_2$ in a 2:1:1 ratio at 15\,K'' ---
measured at a range of temperatures.

The library is not uniform, and normalising it was a substantial
engineering task:

\begin{itemize}
  \item some files store wavenumber (cm$^{-1}$), others wavelength
        ($\mu$m); the loader detects which by checking whether the mean of
        column one exceeds 100, and converts via $\lambda = 10^4/\nu$;
  \item coverage varies --- some files span the whole 5--28\um{} range,
        others only a narrow band around a single absorption feature;
  \item baselines, noise levels and sampling grids differ between files;
  \item about 40 files added from external sources used Windows-1252
        encoding with an en-dash in the header comment, which caused
        \code{np.loadtxt} to fail; the loader was changed to read with
        \code{encoding='latin-1'} and parse line by line.
\end{itemize}

Every loaded spectrum is interpolated onto a fixed grid of 3000 points
spanning 5.0--28.0\um{} and normalised so that its peak value equals 1.
The normalisation is essential: without it, substances with intrinsically
strong absorption would dominate every synthetic mixture regardless of the
concentration actually assigned to them.

\subsection{From raw telescope output to optical depth}

The telescope records flux --- how much light arrives at each wavelength.
The quantity needed for fitting is optical depth, defined as
\[
  \tau(\lambda) \;=\; \ln\!\left(\frac{F_{\text{cont}}(\lambda)}
                                      {F_{\text{obs}}(\lambda)}\right),
\]
where $F_{\text{obs}}$ is the observed flux and $F_{\text{cont}}$ is the
\emph{continuum} --- an estimate of what the flux would have been if no ice
were present. Optical depth is the natural quantity because contributions
from different substances add linearly in $\tau$, which is exactly what
makes the linear model $\mathbf{Ax}\approx\mathbf{b}$ valid.

Two subtractions are therefore needed before any fitting can begin, and
both are error-prone:

\begin{enumerate}
  \item \textbf{Continuum subtraction.} The underlying emission of the
        protostar must be estimated and divided out. Guess it too high and
        spurious absorption appears everywhere; guess it too low and real
        features are erased.
  \item \textbf{Silicate removal.} Rocky dust grains produce a broad
        absorption centred at 9.7\um{} with a secondary feature near
        18\um. This is not ice and must be modelled and removed.
\end{enumerate}

\begin{warnbox}
Errors made at this stage propagate into every downstream result. As
Section~\ref{sec:known-problems} documents, the single largest remaining
limitation of the whole pipeline is a continuum/silicate artifact, not a
failure of the machine learning or the fitting.
\end{warnbox}

\subsection{Three generations of continuum and silicate fitting}

\subsubsection*{Version 1 --- spline continuum (abandoned)}

A \code{UnivariateSpline} was fitted through manually chosen anchor points.
This failed on deeply embedded sources: anchor selection was difficult to
automate, and nothing prevented the spline from dipping into genuine
absorption features and erasing them. Fit quality on the L1448 sources was
only 30--44\%.

\subsubsection*{Version 2 --- log-log polynomial with hard penalty}

This adopted the approach used in Prathap's original INDRA: a degree-4
polynomial in $\log\lambda$ versus $\log F$ space, with one critical
physical constraint built in --- the continuum may never fall below the
observed flux anywhere, because optical depth must be non-negative
everywhere. The constraint was enforced with a hard penalty function and
optimised with Nelder--Mead using hyperparameters $w_{\text{anchor}} =
5.3$, $w_{\text{curve}} = 13.5$.

Fit quality jumped from 30--44\% to 73--74\%. The CH$_4$ band at 7.7\um{}
was recovered (previously zero), and formic acid and formate were detected
for the first time. However, with 34 free parameters the Nelder--Mead
optimiser never converged within its iteration limit on either L1448
source --- it stopped at a local minimum.

\subsubsection*{Version 3 --- L-BFGS-B with a soft barrier (current)}

Three targeted fixes were applied:

\begin{description}[style=nextline, leftmargin=0pt, labelindent=0pt, labelwidth=0pt, labelsep=0pt, itemsep=4pt]
  \item[Fix 1 --- gradient-based optimiser.]
  Nelder--Mead was replaced with L-BFGS-B. Because L-BFGS-B requires
  gradients, the hard penalty function was replaced by a smooth
  exponential barrier of the form $\exp(k \times \text{violation})$, which
  is differentiable everywhere. Convergence now occurs in roughly 1000
  iterations rather than the 200{,}000+ Nelder--Mead required.

  \item[Fix 2 --- additional long-wavelength anchors.]
  Anchor points at 17.5\um{} and 18.0\um{} were added to constrain the
  continuum at the boundary of the water libration region.

  \item[Fix 3 --- two-component silicate model.]
  The single Drude profile was replaced by a physically motivated
  two-component model: olivine (peak 9.8\um) plus pyroxene (peak
  9.65\um), each with a secondary 18\um{} band, using grain parameters
  from Weingartner \& Draine (2001).
\end{description}

\begin{table}[H]
\centering
\small
\begin{tabular}{@{}lccc@{}}
\toprule
\textbf{Metric} & \textbf{V1 (spline)} & \textbf{V2 (Nelder--Mead)} &
\textbf{V3 (L-BFGS-B)} \\
\midrule
IRS2 fit quality        & 30\%  & 74\%          & 74\% (converged) \\
IRS3B fit quality       & 44\%  & 73\%          & 73\% (converged) \\
$\tau$(7.7\um) CH$_4$   & 0.000 & 0.276         & 0.276 \\
Silicate $\tau_{\text{peak}}$ & 0.05 & 8.00 (at bound) & 8.00 (at bound) \\
$\tau$(20\um) artifact  & 0.66 / 0.27 & 1.41 / 1.38 & 1.41 / 1.38 \\
CH$_4$ detected         & No    & Yes           & Yes \\
HCOOH detected          & No    & Yes           & Yes \\
Optimiser converged     & N/A   & No (max iter) & Yes \\
\bottomrule
\end{tabular}
\caption{Continuum fitting across three generations, evaluated on the two
L1448 sources.}
\end{table}

\subsubsection*{The remaining artifact}

Both L1448 sources still show $\tau(20\,\mu\text{m}) \approx 1.4$. This is
physically impossible: no known interstellar ice absorbs beyond 19\um, so
the optical depth there should be essentially zero. The cause is that the
cold envelope emits strongly between 20 and 28\um{} due to thermal dust
emission, and a degree-4 polynomial cannot simultaneously fit the ice
region (5--17\um) and the thermal tail (20--28\um) without undershooting
between 17 and 20\um.

The proposed fix --- not yet implemented --- is to use the physical fact
that $\tau_{\text{ice}}(19\text{--}20\,\mu\text{m}) = 0$ as a hard
calibration anchor for the silicate model's secondary 18\um{} band. Pinning
that band to its correct amplitude should cascade into a corrected silicate
profile across the entire 8--17\um{} region where silicate and ice overlap.

% =====================================================================
\section{Version 3.8 --- The Full Pipeline}
% =====================================================================

This section describes every technique used in v3.8, in the order the data
flows through them.

\subsection{Step 1 --- Spectral clustering of the library}

\subsubsection*{The problem being solved}

Of the 1209 library files, many are temperature variants of the same
mixture --- for instance water:methanol measured at 10, 15, 20, 30, 50,
70, 100, 120 and 150\,K. At the effective resolution of MIRI ($R\sim100$),
several of these are indistinguishable from one another. Training one
classifier per file would mean 1209 classifiers, many of them attempting to
separate signals that are identical within the noise, each trained on very
few examples.

\subsubsection*{The method}

Hierarchical clustering with complete linkage, using cosine distance:
\[
  d_{\cos}(\mathbf{u},\mathbf{v}) \;=\;
  1 - \frac{\mathbf{u}\cdot\mathbf{v}}{\|\mathbf{u}\|\,\|\mathbf{v}\|}.
\]
Cosine distance measures the angle between two spectra viewed as vectors in
3000-dimensional space, ignoring overall scale. This is the correct metric
here: the same ice measured at two different laboratory thicknesses
produces spectra that differ in magnitude but point in the same direction,
and should be treated as identical.

The procedure is:

\begin{enumerate}
  \item Strip the temperature suffix from each filename using the regular
        expression \code{r'\_\textbackslash d+[Kk]\$'} and group the
        variants of each mixture.
  \item Within each group, compute the pairwise cosine distance matrix and
        run complete-linkage hierarchical clustering.
  \item Test 2, 3 and 4 clusters; keep the largest number for which every
        pair of clusters is separated by more than the threshold
        $0.05$. Label the resulting clusters A, B, C, D.
  \item Perform a cross-mixture merge: compute distances between all
        cluster representatives across different mixtures and merge any
        pair closer than $0.05$.
\end{enumerate}

The threshold $0.05$ was chosen to correspond approximately to the MIRI
noise floor: two spectra within 5\% cosine distance cannot be
distinguished at any realistic signal-to-noise ratio.

\begin{goodbox}
\textbf{Result:} 1209 files reduce to \textbf{284 spectral clusters}, a
77\% reduction. Each cluster represents one spectrally distinguishable ice
state and becomes one classification target, with roughly four times as
many training examples per target as before.
\end{goodbox}

A useful validation: when recall was measured both strictly (exact cluster
label) and generously (correct mixture, any temperature class), the gap was
0.0\%. Every miss was a complete mixture miss, never a temperature-class
confusion --- confirming that the clusters are genuinely distinct rather
than arbitrary subdivisions.

\subsection{Step 2 --- Synthetic training data generation}

\subsubsection*{Why synthetic data is necessary}

Supervised learning requires labelled examples: spectra for which the true
composition is known. Real JWST spectra with trusted compositions number in
the single digits. The solution is to generate training data by simulating
the physics.

\subsubsection*{The generator}

135{,}000 synthetic spectra were produced (120{,}000 general plus 15{,}000
targeted). Each is built as follows:

\begin{enumerate}
  \item \textbf{Select components.} Draw a number of layers (10--25,
        peaking at 15) and select that many cluster representatives.
        \htwoo, \cotwo{} and CO are always included as mandatory species.
  \item \textbf{Assign fractions by Dirichlet sampling.} Mixing fractions
        are drawn from a Dirichlet distribution, which guarantees the
        fractions are non-negative and sum to one --- exactly the
        structure of a physical mixture.
  \item \textbf{Scale and sum.} Multiply by a total optical depth and add
        the components linearly.
  \item \textbf{Inject silicate.} In 80\% of spectra, add a synthetic
        silicate feature built from Drude oscillator profiles at 9.7\um{}
        (Si--O stretch) and 18\um{} (O--Si--O bend), with peak optical
        depth sampled from 0.5--8.0. Grain size factor, crystalline
        fraction and olivine/pyroxene ratio are randomised per spectrum.
  \item \textbf{Simulate observation.} Convert to flux, add Gaussian noise
        with $\sigma = 1/\text{SNR}$ for SNR drawn uniformly from
        $U(50,200)$, apply median smoothing and a randomised polynomial
        continuum, and convert back to optical depth.
\end{enumerate}

\subsubsection*{Three generator bugs that had to be found}

The generator went through several revisions, and three bugs in earlier
versions are worth recording because each silently degraded model quality:

\begin{description}[style=nextline, leftmargin=0pt, labelindent=0pt, labelwidth=0pt, labelsep=0pt, itemsep=4pt]
  \item[Temperature normalisation.] Temperatures were divided by 200\,K
        while the true maximum in the database is 150\,K, placing some
        normalised values above 1.0. Fixed by dividing by 150.

  \item[Noise too small.] Noise $\sigma$ was hard-coded at $10^{-4}$,
        roughly 50--100 times smaller than real MIRI noise. The model
        trained on unrealistically clean data and failed to transfer.
        Fixed by sampling SNR per spectrum.

  \item[No normalisation before mixing.] Laboratory spectra were summed
        without first normalising each to peak 1, so substances with
        intrinsically large absorption dominated regardless of their
        assigned Dirichlet fraction. Fixed by normalising at load time.
\end{description}

\subsubsection*{Targeted generation for weak species}

Some species are genuinely hard to learn because their absorption features
are intrinsically weak. If such a species appears in training at a typical
5\% mixing fraction, its contribution may fall below the noise floor --- so
the model receives a positive label with no detectable signal in the input,
and correctly learns to predict near-zero probability.

Two mechanisms address this:

\begin{itemize}
  \item A list of \code{HARD\_SPECIES\_KW} (SO$_2$, H$_2$CO, CH$_3$OCHO,
        CH$_3$CH$_2$OH, \ocn, CH$_3$COOH, HCOOH, CH$_3$CHO and
        variants). Any species on this list receives a Dirichlet
        concentration parameter of $\alpha = \max(\alpha_{\text{current}},
        5.0)$, forcing its features above the noise floor during training.
  \item 15{,}000 targeted spectra in which one critical species is forced
        to appear prominently ($\alpha = 8$) alongside dominant water
        ($\alpha = 10$); and in 40\% of general spectra, one critical
        species is guaranteed to co-occur at $\alpha = 4$.
\end{itemize}

\subsection{Step 3 --- The LightGBM classifiers}

\subsubsection*{Choice of model and metric}

The modelling approach evolved through several stages before settling:
Random Forest regression $\rightarrow$ XGBoost $\rightarrow$ LightGBM
binary classification. Two decisions drove this.

First, \textbf{classification rather than regression}. Predicting the
fractional abundance of each species directly is numerically unstable when
96\% of all labels are zero. Predicting presence or absence, and leaving
quantification to Stage 2, is far better conditioned.

Second, \textbf{PR-AUC rather than $R^2$ or F1} as the primary metric.
$R^2$ is meaningless at 96\% sparsity --- predicting all zeros scores well.
F1 depends on an arbitrary threshold choice. Precision--Recall AUC
summarises the entire precision/recall trade-off without committing to a
threshold, and is the standard choice for heavily imbalanced detection
problems.

\subsubsection*{Recall-first philosophy}

One binary classifier is trained per cluster, each answering ``is this
spectral cluster present?'' from the full 3000-point optical depth
spectrum. Decision thresholds are tuned to a 90\% recall target rather than
to maximise precision.

This asymmetry is deliberate and is the architectural reason the two-stage
design works:

\begin{keybox}
A missed detection in Stage 1 is unrecoverable --- a species never offered
to the fitting stage can never be found. A false alarm in Stage 1 is
cheap, because Stage 2 fits actual spectral shapes and the non-negativity
constraint drives genuinely absent species to exactly zero. The classifiers
are therefore tuned to over-report, and the fitting stage acts as the
filter.
\end{keybox}

\subsubsection*{Two-tier training}

Species on the hard list receive stronger classifiers:

\begin{table}[H]
\centering
\small
\begin{tabular}{@{}lcc@{}}
\toprule
\textbf{Hyperparameter} & \textbf{Hard species} & \textbf{Normal species} \\
\midrule
\code{n\_estimators}      & 600  & 400  \\
\code{learning\_rate}     & 0.03 & 0.05 \\
\code{max\_depth}         & 7    & 6    \\
\code{colsample\_bytree}  & 0.7  & 0.6  \\
\code{min\_child\_samples}& 5    & 10   \\
\code{pos\_weight} cap    & 80   & 40   \\
\bottomrule
\end{tabular}
\caption{Two-tier LightGBM configuration.}
\end{table}

Class imbalance is handled per species rather than globally:
\[
  \texttt{pos\_weight} \;=\;
  \operatorname{clip}\!\left(\frac{n_{\text{neg}}}{n_{\text{pos}}}
  \times 0.5,\; 5.0,\; \text{cap}\right).
\]
A species present in 8\% of spectra receives a weight near 6; one present
in 1\% receives the cap. Previously a single fixed weight of 15 was applied
to all 284 classifiers regardless of their actual prevalence.

\subsubsection*{A memory bug worth recording}

An early training script used \code{joblib} to spawn eight worker
processes, each holding its own 1.3\,GB copy of the training matrix ---
10.4\,GB total, which overflowed available RAM and silently fell back to
slow threading. The fix was to run the outer loop sequentially and let
LightGBM use its own internal parallelism, keeping a single copy of the
training data in memory. Training takes approximately 8--10 hours.

\subsubsection*{The double-tau transform bug}

The most consequential bug in the project's history. The generator saves
its output already converted to optical depth (range 0--6.2). An earlier
training script applied the conversion \emph{again}:
\code{-log(clip(X, 0.002, 1.05))}. This inverted the representation ---
strong absorption became near-zero and continuum became 6.2 --- so the
model was trained on an inverted signal while prediction fed it correctly
oriented data. The fix is a single line, \code{X\_tau = np.maximum(X\_raw,
0.0)}, and it alone raised PR-AUC by nearly 10 percentage points and
precision at the standard threshold from 29.3\% to 57.4\%.

\subsubsection*{Patch classifiers for \ocn}

Two additional classifiers were trained separately for the cyanate ion
(\ocn), using spectral-correlation labelling. This species was
underrepresented in the original library and its primary absorption at
4.62\um{} falls just outside the MIRI range, making it hard for the main
classifiers to detect. The patch classifiers achieve PR-AUC of 94.6\% and
98.7\%, bringing the deployed total to 286.

A note on chemistry that matters for the bookkeeping: \ocn{} at 80\,K is
the thermal reaction product of HNCO + NH$_3$ $\rightarrow$ \nhfour{} +
\ocn. Because the HNCO fundamental lies outside the MIRI range anyway, the
pipeline uses \ocn{} detection to credit both parent species, encoded in a
\code{PARENT\_RULES} table.

\subsection{Step 4 --- The NNLS quantification stage}

With a candidate list from Stage 1, the fitting stage must determine
abundances. Four techniques combine here.

\subsubsection*{Candidate selection by score floor}

Earlier versions took a fixed number of top-scoring candidates
(\code{N\_CANDIDATES = 50} or 120). A sensitivity analysis showed a clear
plateau:

\begin{table}[H]
\centering
\small
\begin{tabular}{@{}lccccc@{}}
\toprule
Candidates $N$ & 20 & 30 & 40 & 50 & 120 \\
\midrule
Recall & 11/15 & 12/15 & 12/15 & \textbf{14/15} & 14/15 \\
\bottomrule
\end{tabular}
\caption{Candidate count sensitivity. Ranks 51--120 add no recall and only
feed false positives into the solver.}
\end{table}

However, a fixed $N$ tuned on one source will not transfer. It was replaced
by a score floor, \code{ML\_SCORE\_FLOOR = 0.05}: take every species
scoring above 5\%. This self-scales --- a chemically simple source produces
fewer candidates, a complex one more.

\subsubsection*{Deduplication}

Candidate spectra within cosine distance 0.015 of an already-selected
candidate are dropped, to avoid feeding near-collinear columns to the
solver.

\subsubsection*{Tikhonov regularisation}

Plain NNLS tends to concentrate all weight on a few dominant species
(water and \cotwo) and zero out everything else. Adding a penalty term,
\[
  \min_{\mathbf{x}\ge 0}\;
  \|\mathbf{A}\mathbf{x}-\mathbf{b}\|^2 + \alpha\|\mathbf{x}\|^2 ,
\]
discourages this concentration. It is implemented by solving ordinary NNLS
on the augmented system
$\left[\begin{smallmatrix}\mathbf{A}\\ \sqrt{\alpha}\,\mathbf{I}
\end{smallmatrix}\right]\mathbf{x} =
\left[\begin{smallmatrix}\mathbf{b}\\ \mathbf{0}\end{smallmatrix}\right]$,
with $\alpha = 0.15$.

\subsubsection*{Orthogonal Matching Pursuit (OMP)}

Rather than solving for all candidates simultaneously, OMP builds the
solution greedily:

\begin{enumerate}
  \item Start with an empty support set and residual $\mathbf{r} =
        \mathbf{b}$.
  \item Find the column of $\mathbf{A}$ most positively correlated with
        $\mathbf{r}$.
  \item Add it to the support and refit NNLS on the support only.
  \item Repeat until the best correlation falls below 0.05 or the maximum
        of 25 components is reached.
\end{enumerate}

This produces sparse solutions (10--25 components) that are faster to
compute and physically more interpretable. When several columns have
similar correlation --- within a 5\% band of the leader --- the ML scores
act as a tiebreaker, preferring the species the classifier is more
confident about.

\subsubsection*{Forced-floor fitting}

Some species are detected confidently by the classifier but driven to
exactly zero by the solver, because a dominant neighbour explains their
absorption region adequately. To prevent this, such species are given a
lower bound of 0.002 rather than 0, and the problem is solved as a bounded
least-squares problem using \code{lsq\_linear} with the BVLS method. This
directly rescued CH$_3$COOH, HCOOH and \nhfour{} in testing.

\subsubsection*{Backward elimination}

After OMP and forcing, 25--33 components typically survive --- more than
physically plausible. Backward elimination prunes them: the
lowest-abundance component is tested for removal, and if removing it does
not degrade the fit, it is dropped. This repeats until no further removal
is safe. There is no fixed target count; the procedure stops naturally,
at 8 components for a simple source and at 20 for a complex one.

\subsection{Version 3.8 results on IRAS 2A}

\begin{table}[H]
\centering
\small
\begin{tabular}{@{}lccccl@{}}
\toprule
\textbf{Version} & \textbf{Fit} & \textbf{ML recall} &
\textbf{NNLS recall} & \textbf{Survivors} & \textbf{Key change} \\
\midrule
v3.7 baseline        & 92.1\% & 27\%  & 53\%  & 23 & Initial pipeline \\
v3.8 OMP + floor     & 92.3\% & 87\%  & 87\%  & 33 & OMP, forced floor \\
v3.8 backward elim.  & 90.6\% & 93\%  & 87\%  & 15 & Pruning added \\
v3.8patch + \ocn     & 90.6\% & 100\% & 100\% & 16 & \ocn{} classifiers \\
\bottomrule
\end{tabular}
\caption{Version 3.8 progression on the benchmark source. The final row
reports perfect recall --- a result that Section~\ref{sec:leakage} shows
cannot be trusted.}
\end{table}

% =====================================================================
\section{The Data Leakage Problem}
\label{sec:leakage}
% =====================================================================

\subsection{What data leakage means}

Data leakage occurs when information that would not be available at
prediction time influences the prediction. The classic textbook example is
accidentally including the target variable among the input features: the
model appears to perform superbly, and the performance is entirely
illusory, because the information it depends on will not exist when the
model is deployed.

Leakage is dangerous precisely because it does not look like a bug. The
code runs, the metrics improve, and the improvement is reported as
progress. It is detected by reasoning about information flow, not by
testing.

\subsection{The leakage in v3.8}

\begin{warnbox}
\textbf{Critical finding.} The v3.8patch script that reported 15/15 recall
was consulting the published ground truth during prediction. It was
forcing components into the solution \emph{because} they appeared in the
published answer, not because the spectral evidence supported them. The
reported figure measures the pipeline's ability to reproduce an answer it
had already been shown.
\end{warnbox}

Two separate mechanisms were identified.

\subsubsection*{Mechanism 1 --- \code{TARGET\_SPECIES} forcing}

The function \code{force\_missing\_components()} compared the current
solution against a hard-coded dictionary of the species known to be present
in IRAS 2A, identified which were absent from the solution, and then
searched the library specifically for files that would fill those gaps.

This is equivalent to a student checking their answers against the marking
scheme mid-examination and correcting the ones that disagree. Whatever the
resulting score, it does not measure what the examination was designed to
measure.

\subsubsection*{Mechanism 2 --- recall guard in backward elimination}

The pruning step removed components one at a time, but before each removal
it checked whether removing that component would reduce recall measured
against the published list. Components were protected from removal
specifically to preserve that recall number.

This is the subtler of the two, and arguably the more instructive. The
intention was reasonable --- ``do not throw away something we know is
correct'' --- but the stopping criterion for an algorithm must never depend
on the answer the algorithm is meant to produce. On a new source with no
published list, this criterion simply does not exist, so the pipeline would
behave differently there in an unpredictable way.

\subsection{Why this mattered more than usual}

In most machine learning projects leakage inflates a reported metric. Here
the consequence was sharper, because the entire purpose of the project is
generalisation to sources with no ground truth. A pipeline tuned --- even
implicitly --- to reproduce the one source with a published answer offers
no evidence whatsoever about its behaviour on L1448 IRS2 or IRS3B.

Several parameter choices were re-examined in this light and found to have
been set, knowingly or not, by what worked on the benchmark:

\begin{itemize}
  \item \code{N\_CANDIDATES = 50} and \code{TARGET\_SURVIVORS = 15} were
        selected because they produced the best result on IRAS 2A;
  \item a forcing threshold of 0.15 had been chosen because it happened to
        rescue the specific species IRAS 2A was missing;
  \item ``recall-aware forcing'' --- forcing a species if doing so
        increased recall --- was leakage stated directly as a design
        feature.
\end{itemize}

\subsection{The structural fix}

The repair was not to remove individual lines of offending code, but to
change the structure of the program so that the error cannot recur:

\begin{goodbox}
\textbf{The published ground truth (\code{paper\_mixtures}) is confined to
a single reporting function that executes after every algorithmic decision
has been made.} It can read the final result; it cannot influence it. No
function involved in selection, forcing, protection or elimination accepts
it as an argument.

\vspace{4pt}
This is enforced by the call structure, not by discipline. Reintroducing
the leak would require deliberately threading ground truth through function
signatures that do not currently accept it.
\end{goodbox}

Equally important, the stopping criterion for backward elimination was
changed from a recall criterion to a \emph{fit-quality} criterion: a
component is removed if removing it degrades the residual by less than
0.5\%. Fit quality is computable on any source, with or without a published
answer, so the algorithm now behaves identically on the benchmark and on
new sources.

% =====================================================================
\section{Version 3.9 --- The Clean Pipeline}
% =====================================================================

Version 3.9 is a rewrite of the prediction and fitting stages. The
classifiers themselves were not retrained --- the leakage was entirely
downstream of Stage 1.

\subsection{Architecture}

\begin{keybox}
\textbf{Stage 1 --- Detection.} LightGBM classifiers score all 286 spectral
clusters against the \emph{silicate-included} optical depth spectrum.
Species clearing the score floor become candidates.

\vspace{6pt}
\textbf{Stage 2 --- Quantification.} Multi-component NNLS, initialised by
OMP, fits the \emph{silicate-removed} ice-only spectrum using the candidate
laboratory spectra as basis vectors.
\end{keybox}

\subsubsection*{Why two different spectra}

This asymmetry is deliberate and frequently misunderstood. The classifiers
were trained on synthetic spectra that \emph{include} silicate absorption,
so at inference they must be shown a spectrum with silicate present. The
laboratory reference spectra used for fitting contain \emph{no} silicate,
so the target for fitting must have silicate removed. Using the wrong
spectrum for either stage degrades performance substantially.

Where only the silicate-removed spectrum is available --- as is the case
for the IRAS 2A files --- silicate is reconstructed with a Drude profile
and added back before classification.

\subsection{Option A --- ML-score parent forcing (primary method)}

After OMP selects an initial component set, the pipeline checks whether
every parent species with an ML score of at least 20\% is represented in
the solution. Any that is missing is forced in at a floor abundance of
0.002, using its highest-scoring candidate file.

The critical property is that this rule consults \emph{only classifier
output}. It makes no reference to any published result, so it executes
identically on IRAS 2A and on a source nobody has ever analysed. The
threshold of 20\% is fixed and generic, deliberately replacing the 0.15
that had been tuned to the benchmark.

\subsection{Option C --- confidence-weighted Tikhonov NNLS}

An alternative formulation replaces uniform regularisation with per-column
weights derived from classifier confidence:
\[
  \alpha_j \;=\; \frac{\alpha}{s_j + \varepsilon},
\]
where $s_j$ is the ML score of column $j$. High confidence yields a low
penalty and permits a large coefficient; low confidence yields a high
penalty and pushes the coefficient toward zero. This suppresses
low-confidence false positives smoothly, with no hard thresholds anywhere.

It is elegant, and it underperforms: 12/15 versus 15/15. Without explicit
forcing, minor species present at genuinely small abundances are penalised
to zero along with the false positives. Option A remains the primary
method; Option C is retained as a comparison that demonstrates the forcing
step is doing real work rather than merely patching over a weakness.

\subsection{OMP protection}

Any component that OMP selected independently \emph{and} that carries an ML
score of at least 10\% is protected from backward elimination. The
reasoning is that two independent lines of evidence support it: OMP chose
it on spectral residual grounds, and the classifier assigned it meaningful
confidence on pattern-recognition grounds.

This threshold is flagged internally as probably too permissive --- it
weakens elimination more than intended --- so every protected component is
tagged in the output rather than silently accepted.

\subsection{Diagnostic flags}

Three flags are attached to every reported component, so that weaknesses
appear in the output rather than hiding inside it:

\begin{description}[style=nextline, leftmargin=0pt, labelindent=0pt, labelwidth=0pt, labelsep=0pt, itemsep=4pt]
  \item[\code{protected\_only}] The component survived only because
        protection shielded it from elimination, not because independent
        fitting confirmed it.
  \item[\code{solver\_override\_suspect}] The solver assigned a
        substantial abundance to a component the classifier scored low.
        This usually indicates the solver is absorbing a band-shape
        artifact rather than detecting the species.
  \item[Quantitative / probabilistic split] Components with fitted
        fraction above 1\% are reported as quantitative detections; those
        below are reported as probable detections at upper-limit abundance.
\end{description}

\subsection{The honest assessment of 15/15}

Version 3.9 Option A achieves 15/15 recall on IRAS 2A with no leakage, at
91.1\% fit quality and 21 surviving components. That number requires
careful statement:

\begin{table}[H]
\centering
\small
\begin{tabular}{@{}lll@{}}
\toprule
\textbf{Category} & \textbf{Count} & \textbf{Basis} \\
\midrule
Quantitative detections & $\sim$13 &
Fitted fraction $>1\%$, confirmed by both OMP and classifier \\
Floor-level detections  & $\sim$6 &
Fraction $<0.5\%$, present only because forcing inserted them \\
\bottomrule
\end{tabular}
\caption{Breakdown of the 15/15 result.}
\end{table}

\begin{goodbox}
\textbf{The defensible claim:} 13 species detected quantitatively, plus 6
flagged as probable detections at upper-limit abundance. The pipeline
reached 15/15 recall without using any source-specific information --- but
6 of those 15 are detection claims rather than abundance measurements, and
must be reported as upper limits.
\end{goodbox}

This is a weaker claim than ``15/15 recall'' and a considerably stronger
one than the leaked v3.8 figure, because it is accompanied by an explicit
account of what each number does and does not establish.

% =====================================================================
\section{Results}
% =====================================================================

\subsection{Benchmark validation --- IRAS 2A}

\begin{table}[H]
\centering
\small
\begin{tabular}{@{}lccccl@{}}
\toprule
\textbf{Version} & \textbf{Fit} & \textbf{ML recall} &
\textbf{NNLS recall} & \textbf{Survivors} & \textbf{Status} \\
\midrule
v3.7 baseline       & 92.1\% & 27\%  & 53\%   & 23 & --- \\
v3.8 OMP + floor    & 92.3\% & 87\%  & 87\%   & 33 & --- \\
v3.8 backward elim. & 90.6\% & 93\%  & 87\%   & 15 & --- \\
v3.8patch + \ocn    & 90.6\% & 100\% & 100\%  & 16 & \textcolor{warnred}{\textbf{leaked}} \\
v3.9 Option A       & 91.1\% & ---   & 15/15  & 21 & \textcolor{goodgreen}{\textbf{clean}} \\
v3.9 Option C       & 91.1\% & ---   & 12/15  & 17 & \textcolor{goodgreen}{\textbf{clean}} \\
\bottomrule
\end{tabular}
\caption{Full version history on the benchmark source.}
\end{table}

\subsection{Blind generalisation --- L1448 IRS2 and IRS3B}

Neither source has a published decomposition, so there is no recall figure
to report. The table below \emph{is} the result, not a comparison against
anything.

\begin{table}[H]
\centering
\small
\begin{tabular}{@{}lcc@{}}
\toprule
\textbf{Property} & \textbf{L1448 IRS2} & \textbf{L1448 IRS3B} \\
\midrule
Type                        & Class 0 protostar & Class 0 binary \\
Luminosity                  & 5.2\,$L_\odot$    & Unknown (binary) \\
Distance                    & $\sim$300\,pc     & $\sim$300\,pc \\
$\tau$(6\um) \htwoo         & 1.883             & 1.872 \\
$\tau$(15.2\um) \cotwo      & 3.170             & 2.493 \\
$\tau$(20\um) artifact      & 1.410             & 1.379 \\
NNLS fit quality            & 73.9\%            & 73.4\% \\
Surviving components        & 13                & 13 \\
\bottomrule
\end{tabular}
\caption{Blind run results on the two new sources.}
\end{table}

\subsubsection*{Species detected in both sources}

\htwoo{} (dominant), \cotwo, CO, CH$_4$, NH$_3$, \nhfour, H$_2$S, HCOOH
(formic acid), HCOO$^-$ (formate), \ocn{} and organic residue. The CH$_4$
detection at $\tau(7.7\,\mu\text{m}) = 0.28$ and $0.40$ is the first for
these sources.

That both sources return the same species set and the same survivor count
is meaningful evidence: it suggests the pipeline is responding to
physically consistent ice chemistry rather than to source-specific noise.

\subsubsection*{Species absent, and why}

\begin{goldbox}
\textbf{\chthreeoh{} (methanol) is not detected in either source.} Its
primary C--O stretch at 9.74\um{} falls exactly where silicate
over-subtraction drives the ice-only optical depth to zero. This is a
continuum artifact, not a genuine non-detection, and it is the single
clearest demonstration that the remaining continuum problem has real
scientific cost.
\end{goldbox}

SO$_2$ and H$_2$CO are absent from L1448 but present in IRAS 2A. This is
physically reasonable rather than suspicious: IRAS 2A is roughly seven
times more luminous, producing a warmer inner envelope with more complex
chemistry, whereas L1448 IRS2 at 5.2\,$L_\odot$ would be expected to have a
simpler ice inventory.

The lower fit quality on the new sources (73--74\% versus 91\%) traces
directly to the long-wavelength continuum artifact, not to the
classification or fitting logic.

% =====================================================================
\section{Library Improvement --- Gaussian Peak Decomposition}
% =====================================================================

A substantial side effort replaced the raw laboratory files with smooth
analytic representations. Each spectrum is decomposed into a sum of
asymmetric Gaussian profiles, one per absorption band.

\subsection{Motivation}

\begin{itemize}
  \item \textbf{Clean basis vectors.} Raw files carry varying baselines,
        noise spikes and irregular sampling grids; fitted Gaussians are
        smooth and artifact-free.
  \item \textbf{Temperature interpolation.} Once peak position, width and
        amplitude are fitted at the discrete measured temperatures, they
        can be interpolated continuously, allowing the fitting stage to
        work at any temperature rather than only the measured ones.
  \item \textbf{Controllable band ratios.} Relative band amplitudes within
        a mixture become explicit parameters during synthetic data
        generation.
  \item \textbf{Better separation of weak species.} A cleaner basis helps
        distinguish SO$_2$ at 8.7\um{} and CH$_3$CH$_2$OH at 9.5\um{} from
        the surrounding silicate absorption.
\end{itemize}

\subsection{Implementation}

\begin{itemize}
  \item \textbf{Baseline subtraction:} Asymmetric Least Squares (ALS) from
        the \code{pybaselines} library.
  \item \textbf{Peak finding:} \code{scipy.signal.find\_peaks} with
        prominence and width constraints tuned for ice bands (minimum
        prominence 4\% of peak, minimum width 5\,cm$^{-1}$).
  \item \textbf{Fitting:} LMFIT with Levenberg--Marquardt optimisation.
        Each peak uses a different $\sigma$ on its left and right sides, to
        capture the characteristic asymmetry of phonon-broadened ice bands.
  \item \textbf{Parallel execution:} \code{ProcessPoolExecutor} with seven
        workers.
\end{itemize}

\begin{table}[H]
\centering
\small
\begin{tabular}{@{}ll@{}}
\toprule
\textbf{Result} & \textbf{Value} \\
\midrule
Files processed successfully & 1198 / 1209 \\
Files failed                 & 12 (11 with no peaks, 1 zero after baseline) \\
Processing time              & 15\,min 55\,s (7 workers, vs.\ 40+ min serial) \\
Mean peaks per file          & 7.5 \\
Mean fit RMS                 & 0.381 \\
\bottomrule
\end{tabular}
\caption{Batch decomposition over the full library.}
\end{table}

\begin{goldbox}
\textbf{Known issue.} The mean RMS of 0.381 is high. Two causes: mixture
files with broad overlapping bands need more than 7--8 Gaussian components,
and the ALS baseline under-performs on high-absorbance files ($\tau > 3$)
because its smoothing parameter $\lambda$ is not yet adaptive. This must be
fixed before the Gaussian library can be used to retrain the classifiers.
\end{goldbox}

An interactive visualiser (\code{gaussian\_visualizer.py}) was also built,
displaying the raw spectrum with its fitted baseline, the individual
Gaussian components, the residuals and the reconstruction, for inspection
of any file on demand.

% =====================================================================
\section{Current State and Known Limitations}
\label{sec:known-problems}
% =====================================================================

\subsection{Working reliably}

\begin{table}[H]
\centering
\small
\begin{tabular}{@{}L{5.2cm}L{9.6cm}@{}}
\toprule
\textbf{Component} & \textbf{Notes} \\
\midrule
LightGBM classifiers (286) & Trained and validated; not yet retrained on
Gaussian data \\
\ocn{} patch classifiers (2) & PR-AUC 94.6\% and 98.7\% \\
Continuum fitting (L-BFGS-B) & Converges reliably on all three sources \\
Two-component silicate model & Working, but reaches its upper bound on
L1448 \\
OMP + backward elimination & Clean, source-agnostic, no leakage \\
Option A parent forcing & Deployed, generic 20\% threshold \\
Gaussian peak database & 1198/1209 files; baseline fix needed \\
IRAS 2A validation & 91.1\% fit, 13 quantitative + 6 floor detections \\
L1448 IRS2 / IRS3B & Blind run complete, 13 species each \\
\bottomrule
\end{tabular}
\end{table}

\subsection{Known problems}

\begin{description}[style=nextline, leftmargin=0pt, labelindent=0pt, labelwidth=0pt, labelsep=0pt, itemsep=5pt]
  \item[Long-wavelength optical depth artifact.]
  $\tau(20\,\mu\text{m})\approx1.4$ on both L1448 sources, where physics
  requires approximately zero. Root cause: a degree-4 polynomial cannot fit
  both the ice region and the thermal tail simultaneously.

  \item[Silicate over-subtraction.]
  The silicate peak optical depth reaches its upper bound of 8.0 on the
  L1448 sources, zeroing the 9.7\um{} region. This directly prevents
  methanol detection (C--O stretch at 9.74\um) and ammonia detection
  (umbrella mode at 9.34\um). Same root cause as above.

  \item[Gaussian baseline quality.]
  ALS fails on high-absorbance mixture files; adaptive $\lambda$ selection
  is required before the Gaussian library can be used for retraining.

  \item[Floor-level detections.]
  Six of the fifteen species detected on IRAS 2A sit at forced floor
  abundance and are not independently quantified. They must be reported as
  upper limits.

  \item[Suspicious \cotwo{} isotopologue component.]
  One component contributes 14.6\% of the total fitted abundance at only
  5.1\% classifier confidence. This pattern --- large coefficient, low
  confidence --- suggests the solver is absorbing a \cotwo{} band-shape
  mismatch rather than genuinely detecting the isotopologue. It is flagged
  by \code{solver\_override\_suspect}.

  \item[Option C underperformance.]
  Confidence-weighted NNLS gives 12/15 rather than 15/15, because it has no
  explicit forcing for minor species at low classifier confidence.
\end{description}

% =====================================================================
\section{Roadmap}
% =====================================================================

\begin{table}[H]
\centering
\small
\begin{tabular}{@{}L{2.4cm}L{4.8cm}L{4.9cm}L{2.3cm}@{}}
\toprule
\textbf{Priority} & \textbf{Upgrade} & \textbf{Expected impact} &
\textbf{Status} \\
\midrule
1 --- Critical & Silicate continuum with $\tau = 0$ constraint beyond
19\um & Fixes methanol detection; removes long-wave artifact & Planned \\
2 --- High & Adaptive ALS $\lambda$ for Gaussian baselines & Enables
reliable Gaussian library for retraining & In progress \\
3 --- High & Retrain classifiers on Gaussian-derived synthetic data &
Improves SO$_2$, H$_2$CO, CH$_3$CH$_2$OH detection & Pending \#2 \\
4 --- Medium & Noise-weighted NNLS ($1/\sigma_i$ weights) & $+3$--5\% fit
quality & Planned v3.10 \\
5 --- Medium & Use the curated baseline-corrected library as primary basis
& $+3$--4\% fit quality & Planned v3.11 \\
6 --- Medium & \code{optool} synthetic silicate profiles & Resolves the
silicate upper-bound problem & Planned v3.12 \\
7 --- Lower & Run on IRAS 4A, B335 and other Class 0 sources &
Generalisation validation beyond Perseus & Pending data \\
\bottomrule
\end{tabular}
\end{table}

\subsection{What can and cannot be claimed today}

\begin{goodbox}
\textbf{Defensible now:}
\begin{itemize}
  \item 15/15 recall on IRAS 2A without using published ground truth
        during prediction.
  \item 91\% fit quality on the benchmark source.
  \item Physical detections of \htwoo, \cotwo, CO, CH$_4$, NH$_3$,
        \nhfour, H$_2$S, HCOOH, HCOO$^-$, \ocn{} and organic residue in
        both L1448 sources.
  \item A two-stage architecture that runs on a new Class 0 protostar with
        no per-source tuning.
\end{itemize}
\end{goodbox}

\begin{goldbox}
\textbf{Needs strengthening before a science paper:}
\begin{itemize}
  \item Methanol detection requires fixing the silicate subtraction.
  \item Weak species (SO$_2$, H$_2$CO, CH$_3$CH$_2$OH) need classifiers
        retrained on Gaussian-derived data.
  \item Fit quality on new sources must reach $\ge 85\%$ (currently
        73--74\%).
  \item At least three to four additional sources are needed before a
        generalisation claim is supportable.
\end{itemize}
\end{goldbox}

% =====================================================================
\section{Summary of Techniques Applied}
% =====================================================================

For quick reference, every substantive technique in the pipeline and the
reason it was chosen:

\small
\begin{longtable}{@{}L{5.0cm}L{10.0cm}@{}}
\toprule
\textbf{Technique} & \textbf{Why it was used} \\
\midrule
\endfirsthead
\toprule
\textbf{Technique} & \textbf{Why it was used} \\
\midrule
\endhead
Hierarchical clustering, complete linkage, cosine distance &
Merge spectrally indistinguishable temperature variants; 1209 files
$\rightarrow$ 284 targets \\
Dirichlet mixture sampling &
Generates non-negative fractions summing to one --- the natural model of a
physical mixture \\
Drude oscillator silicate model &
Physically motivated profile for mineral dust absorption \\
Synthetic data generation (135k spectra) &
Supervised learning requires labels; real labelled spectra are almost
nonexistent \\
LightGBM binary classification &
Fast, handles class imbalance well; binary framing avoids the instability
of regression at 96\% sparsity \\
PR-AUC as primary metric &
$R^2$ is meaningless at high sparsity; F1 depends on an arbitrary threshold \\
Recall-first threshold tuning (90\%) &
A missed detection is unrecoverable; a false alarm is removed later by the
fitting stage \\
Per-species \code{scale\_pos\_weight} &
Species prevalence ranges from 1\% to 8\%; one global weight cannot serve
both \\
Two-tier hyperparameters &
Intrinsically weak species need more capacity and higher positive weight \\
Targeted generation, hard-species $\alpha$ boost &
Ensures weak species clear the noise floor in training data \\
Silicate reconstruction before inference &
Classifiers were trained with silicate present and must see it at
prediction time \\
Non-Negative Least Squares &
Mixture coefficients cannot be negative \\
Tikhonov regularisation &
Prevents the solver concentrating all weight on dominant species \\
Orthogonal Matching Pursuit &
Produces sparse, interpretable solutions; avoids the dense solutions of
simultaneous fitting \\
Cosine deduplication (0.015) &
Removes near-collinear columns that destabilise the solver \\
Forced-floor via bounded least squares &
Rescues confidently detected species that the solver zeroes out \\
Backward elimination on fit quality &
Principled pruning with a criterion computable on any source \\
Structural isolation of ground truth &
Makes the data leakage of v3.8 impossible to reintroduce accidentally \\
Quantitative / probabilistic split &
Distinguishes measured abundances from detection claims in the output \\
Asymmetric Gaussian decomposition &
Clean, interpolatable basis vectors to replace noisy raw files \\
L-BFGS-B with exponential barrier &
Differentiable constraint enforcement; converges where Nelder--Mead did not \\
\bottomrule
\end{longtable}
\normalsize

% =====================================================================
\section{Glossary}
% =====================================================================

\small
\begin{longtable}{@{}L{3.6cm}L{11.4cm}@{}}
\toprule
\textbf{Term} & \textbf{Definition} \\
\midrule
\endfirsthead
\toprule
\textbf{Term} & \textbf{Definition} \\
\midrule
\endhead
ALS & Asymmetric Least Squares --- a baseline subtraction algorithm that
fits a smooth curve lying below absorption peaks \\
Class 0 protostar & The youngest stage of star formation, still embedded in
its natal cloud with most mass in the envelope rather than the forming star \\
Cosine distance & $1 - \cos\theta$ between two vectors; measures shape
similarity while ignoring overall scale \\
Drude profile & A mathematical profile from the physics of a damped
harmonic oscillator, used to model silicate dust absorption \\
Fit quality & $1 - \|\mathbf{b}-\mathbf{A}\mathbf{x}\| / \|\mathbf{b}\|$;
how much of the observed signal the fitted model reproduces \\
LIDA & Leiden Ice Database --- 1209 laboratory ice absorption spectra at
various temperatures and mixture compositions \\
LightGBM & A gradient-boosted decision tree framework, faster and more
memory-efficient than standard boosting implementations \\
MIRI-MRS & Mid-InfraRed Instrument, Medium Resolution Spectrograph, on
JWST; covers 5--28\um{} at $R\sim1300$--3700 \\
NNLS & Non-Negative Least Squares --- least squares with the constraint
$\mathbf{x}\ge0$ \\
OMP & Orthogonal Matching Pursuit --- a greedy algorithm that selects
dictionary columns one at a time by residual correlation \\
Optical depth $\tau$ & $\ln(F_{\text{cont}}/F_{\text{obs}})$; larger
$\tau$ means more absorption \\
PR-AUC & Area under the precision--recall curve; the standard metric for
imbalanced detection problems \\
Recall & Fraction of truly present species that the system detects \\
Tikhonov regularisation & An $\|\mathbf{x}\|^2$ penalty added to least
squares to discourage extreme coefficients \\
$\tau$ ice-only & Optical depth with silicate removed; the fitting target,
because laboratory spectra contain no silicate \\
$\tau$ with silicate & Total optical depth including silicate; the
classifier input, because classifiers were trained on such spectra \\
\bottomrule
\end{longtable}
\normalsize

% =====================================================================
\section*{Key References}
\addcontentsline{toc}{section}{Key References}
% =====================================================================

\small
\begin{itemize}[leftmargin=1.2em]
  \item Rayalacheruvu P., Majumdar L., Rocha W.\,R.\,M., et al.\ (2025).
        \textit{ApJS}, \textbf{281}, 51. --- The original INDRA paper;
        IRAS 2A ground truth with 15 components.
  \item Boogert A.\,C.\,A., Pontoppidan K.\,M., Knez C., et al.\ (2008).
        \textit{ApJ}, \textbf{678}, 985. --- Spitzer c2d ice survey;
        Class 0 abundance reference ranges.
  \item Weingartner J.\,C.\ \& Draine B.\,T.\ (2001). \textit{ApJ},
        \textbf{548}, 296. --- Silicate grain composition and optical
        properties used in the two-component Drude model.
  \item Rocha W.\,R.\,M., van Dishoeck E.\,F., Ressler M.\,E., et al.\
        (2024). \textit{A\&A}, \textbf{683}, A124. --- JWST MIRI-MRS
        analysis of IRAS 2A; source of the optical depth files used in
        validation.
  \item Slavicinska K., et al.\ (2025). --- \htwoo:NH$_3$:H$_2$S 10:2:1 at
        135\,K, the highest-scoring classifier file across all three
        sources.
  \item Öberg K.\,I., et al.\ (2007). --- Pure \htwoo{} ice library spectra
        at 15--160\,K, used throughout the fitting basis.
  \item Pontoppidan K.\,M., et al.\ (2008). \textit{ApJ}, \textbf{678},
        1005. --- Pure \cotwo{} grain-shape corrected spectra.
  \item Novozamsky J.\,H., et al.\ (2001). --- HNCO:NH$_3$ at 80\,K, the
        \ocn{} reference spectrum.
\end{itemize}
\normalsize

\vspace{1.2cm}
\begin{tcolorbox}[colback=navy, colframe=navy, arc=2pt,
                  left=12pt, right=12pt, top=10pt, bottom=10pt]
{\color{white}\large\bfseries Bottom line\par}
\vspace{4pt}
{\color{white}
INDRA 2.0 has progressed from a manual, expert-driven tool to an automated
machine learning pipeline that achieves 15/15 recall on its benchmark
source without consulting any source-specific information during
prediction. It has been run blind on two previously unanalysed Class 0
protostars, detecting 13 physically consistent ice species in each. The
principal remaining bottleneck is the accuracy of the silicate subtraction
near 9--10\um, which both prevents methanol detection and inflates the
long-wavelength optical depth. Correcting this, and retraining the
classifiers on Gaussian-derived synthetic data, are the two highest-priority
next steps before the tool is ready for publication.}
\end{tcolorbox}

\end{document}
